% Chapter: Proof Theory Examples
% Part of the CCTT Monograph

\chapter{Proof Theory: Worked Examples}
\label{ch:proof-theory-examples}

This chapter presents complete, worked proofs of algorithmic equivalences
using the CCTT proof theory.  Each example includes the two programs,
the proof strategy, the proof term, and commentary on why the strategy
is appropriate.

\section{Proof Strategies Overview}

Before diving into examples, we summarize the available strategies
and when each is appropriate:

\begin{longtable}{lp{8cm}}
\toprule
\textbf{Strategy} & \textbf{When to Use} \\
\midrule
$\mathsf{Refl}$ & Same algorithm, same code \\
$\mathsf{NatInd}$ & Recursive vs.\ iterative (same recurrence) \\
$\mathsf{ListInd}$ & List-processing equivalences \\
$\mathsf{WFInd}$ & Programs with non-obvious termination \\
$\mathsf{Simulation}$ & Different algorithms, relatable state spaces \\
$\mathsf{AbsRef}$ & Both implement the same abstract spec \\
$\mathsf{CommDiag}$ & Representation change (encoding bijection) \\
$\mathsf{LoopInv}$ & Loop correctness, stability, invariant props \\
$\mathsf{FiberDecomp}$ & Type-dependent behavior (per-fiber proof) \\
$\mathsf{Z3Discharge}$ & Decidable arithmetic/logic subgoals \\
\bottomrule
\end{longtable}


\section{Example 1: Recursive vs.\ Iterative Factorial}

\subsection{The Programs}

\begin{lstlisting}
def fact_rec(n):
    if n == 0: return 1
    return n * fact_rec(n - 1)

def fact_iter(n):
    result = 1
    for i in range(1, n + 1):
        result *= i
    return result
\end{lstlisting}

\subsection{OTerm Representation}

\begin{align*}
f &= \mathsf{fix}[\mathsf{fact\_rec}]\big(
  \mathsf{case}(n = 0,\; 1,\; n \cdot \mathsf{fact\_rec}(n-1))\big) \\
g &= \fold[{\cdot}](1, \mathsf{range}(1, n+1))
\end{align*}

\subsection{Proof Strategy: Natural Number Induction}

\textbf{Property}: $P(n) \equiv (f(n) = g(n))$.

\textbf{Base case} ($n = 0$):
\begin{align*}
f(0) &= \mathsf{case}(\mathsf{true}, 1, \ldots) = 1 \\
g(0) &= \fold[{\cdot}](1, \mathsf{range}(1, 1)) = \fold[{\cdot}](1, []) = 1
\end{align*}

\textbf{Inductive step}: Assuming $f(k) = g(k)$, prove $f(k+1) = g(k+1)$:
\begin{align*}
f(k+1) &= (k+1) \cdot f(k) \quad [\delta\text{-unfold}] \\
        &= (k+1) \cdot g(k) \quad [\text{IH}] \\
        &= (k+1) \cdot \fold[{\cdot}](1, \mathsf{range}(1, k+1)) \\
        &= \fold[{\cdot}](1, \mathsf{range}(1, k+2)) \quad [\text{fold extension}] \\
        &= g(k+1)
\end{align*}

\subsection{Proof Term}
\begin{lstlisting}
proof = NatInduction(
    base_case   = Assume('base_eval', f[n:=0], g[n:=0]),
    ind_step    = Assume('IH_n', f, g),
    variable    = 'n',
    base_value  = OLit(0),
)
\end{lstlisting}


\section{Example 2: Merge Sort vs.\ Quick Sort}

\subsection{The Programs}
Both sort a list of comparable elements.  Merge sort divides in half and
merges; quick sort partitions around a pivot and concatenates.

\subsection{Proof Strategy: Abstraction Refinement}

Both algorithms implement the abstract specification
$\mathsf{sorted\_permutation}$: a function that returns a sorted
permutation of its input.

\begin{itemize}
\item \textbf{Merge sort}: merge preserves elements (permutation); merge
  of two sorted halves is sorted.
\item \textbf{Quick sort}: partition preserves elements; concatenation of
  sorted(${}<{}$pivot) $\mathbin{+\!\!+}$ [pivot] $\mathbin{+\!\!+}$ sorted(${}>{}$pivot) is sorted.
\end{itemize}

Since both are correct implementations of the same spec, they are
equivalent.

\begin{lstlisting}
proof = AbstractionRefinement(
    spec_name     = 'sorted_permutation',
    abstraction_f = Assume('merge_sort_correct', ...),
    abstraction_g = Assume('quick_sort_correct', ...),
)
\end{lstlisting}

\begin{remark}
Abstraction refinement is the go-to strategy when two algorithms solve
the same problem but use entirely different approaches.  The
specification acts as a ``meeting point'' in the space of behaviors.
\end{remark}


\section{Example 3: Edit Distance DP vs.\ Memoized Recursive}

\subsection{The Programs}
The dynamic programming version fills a 2D table bottom-up.
The memoized version computes top-down with caching.

\subsection{Proof Strategy: Simulation}

The simulation relation is:
\[
  R(\mathit{dp}, \mathit{memo}) \equiv
  \forall (i,j) \text{ computed}.\;
  \mathit{dp}[i][j] = \mathit{memo}[(i,j)]
\]

\textbf{Initial states}: Both initialize base cases identically:
$\mathit{dp}[0][j] = j$, $\mathit{memo}[(0,j)] = j$.

\textbf{Step preservation}: Both compute the same recurrence:
\[
  v = \min(\mathit{del} + 1,\; \mathit{ins} + 1,\;
    \mathit{sub} + \mathit{cost})
\]
The DP fills row-by-row; memoization fills on-demand.  Since both
use the same recurrence and all sub-results are available (DP:
filled in order; memo: recursion + cache), the values match.

\textbf{Output}: Both return the value at $(m, n)$.


\section{Example 4: BFS Flood Fill vs.\ Union-Find}

Both count connected components in a grid.

\textbf{Strategy}: Abstraction refinement via
$\mathsf{count\_connected\_components}$.

\begin{itemize}
\item BFS: each BFS traversal from an unvisited cell visits exactly one
  component.  Count of traversals = count of components.
\item Union-Find: union adjacent cells; count distinct roots.  Each root
  represents exactly one component.
\end{itemize}

This is a clean example of abstraction refinement: the two algorithms
have completely different state representations (visited set vs.\ parent
array), but both correctly compute the same abstract quantity.


\section{Example 5: Dijkstra vs.\ Bellman-Ford}

On graphs with non-negative edge weights, both compute single-source
shortest paths.

\textbf{Strategy}: Simulation with distance invariant.
\[
  R \equiv \forall v \text{ finalized}.\;
  \mathit{dist}_{\mathrm{Dijkstra}}[v] = \mathit{dist}_{\mathrm{BF}}[v]
\]

The key insight: for non-negative weights, Dijkstra's greedy choice
(finalize the minimum-distance vertex) produces the same final distances
as Bellman-Ford's exhaustive relaxation.  The simulation relation holds
because:
\begin{enumerate}
\item Both initialize $\mathit{dist}[\mathit{source}] = 0$,
  $\mathit{dist}[v] = \infty$.
\item Relaxation is the same operation: $d[v] \gets \min(d[v], d[u]+w)$.
\item Non-negative weights ensure Dijkstra's greedy order doesn't miss
  shorter paths.
\end{enumerate}


\section{Example 6: Trial Division vs.\ Sieve Factorization}

Both compute the prime factorization of $n$.

\textbf{Strategy}: Abstraction refinement via
$\mathsf{prime\_factorization}$.

By the Fundamental Theorem of Arithmetic, the prime factorization is
unique.  Both algorithms produce it:
\begin{itemize}
\item Trial division: divides by the smallest prime first, extracting
  each factor with multiplicity.
\item Sieve: precomputes the smallest prime factor (SPF) table, then
  repeatedly divides $n$ by $\mathrm{SPF}[n]$.
\end{itemize}


\section{Example 7: Powerset Recursive vs.\ Bitmask}

\textbf{Strategy}: Commutative diagram with bijection.

The recursion tree has $2^n$ paths, each corresponding to a sequence
of include/exclude decisions.  These map bijectively to $n$-bit masks:
\[
\begin{tikzcd}[column sep=4em]
  \text{recursion paths} \arrow[r, "\mathsf{encode}"] \arrow[d, "\mathsf{collect}"]
    & \{0,\ldots,2^n-1\} \arrow[d, "\mathsf{bits\_to\_subset}"] \\
  \mathcal{P}(S) & \mathcal{P}(S) \arrow[l, "="]
\end{tikzcd}
\]

The encoding is an isomorphism: each path maps to a unique bitmask
and vice versa.  The commutative diagram establishes that both
enumeration methods produce the same power set.


\section{Example 8: LCS via DP vs.\ Hunt-Szymanski}

The standard DP fills a 2D table; Hunt-Szymanski reduces LCS to
Longest Increasing Subsequence (LIS) of match positions.

\textbf{Strategy}: Simulation relating DP entries to LIS lengths.

The key insight: $\mathit{dp}[i][j]$ equals the length of the LIS of
match positions up to $(i,j)$.  Both algorithms compute the same
recurrence; they differ only in evaluation order and pruning strategy.


\section{Example 9: Red-Black Tree vs.\ AVL Tree Insert}

Both are balanced BSTs that support the same interface.

\textbf{Strategy}: Abstraction refinement via BST in-order invariant.

After any sequence of insertions, the in-order traversal of both trees
is the same sorted sequence.  The balancing strategies differ
(red-black rotations vs.\ AVL height-balance rotations), but both
maintain the BST property and produce the same sorted order.


\section{Example 10: Regex NFA vs.\ DFA Matching}

\textbf{Strategy}: Simulation via subset construction.

The simulation relation is:
\[
  R \equiv \mathit{dfa\_state} = \mathsf{frozenset}(\mathit{nfa\_states})
\]

This is the classical subset construction bisimulation.  At each input
character, the DFA transitions to the set of all NFA states reachable
from the current set, which is exactly the set of active NFA states.


\section{Example 11: Map-then-Filter vs.\ Filter-then-Map}

\textbf{Strategy}: List induction with equational reasoning.

For pure $f$ and predicate $p$:
\[
  \mathsf{filter}(p, \mathsf{map}(f, \mathit{xs}))
  = \mathsf{map}(f, \mathsf{filter}(p \circ f, \mathit{xs}))
\]

\textbf{Base case}: both sides on $[]$ give $[]$.

\textbf{Cons case}: assuming the equality for $\mathit{xs}$,
\begin{align*}
  &\mathsf{filter}(p, \mathsf{map}(f, x::\mathit{xs})) \\
  &= \mathsf{filter}(p, f(x) :: \mathsf{map}(f, \mathit{xs})) \\
  &= \begin{cases}
    f(x) :: \mathsf{filter}(p, \mathsf{map}(f, \mathit{xs}))
      & \text{if } p(f(x)) \\
    \mathsf{filter}(p, \mathsf{map}(f, \mathit{xs}))
      & \text{otherwise}
  \end{cases} \\
  &= \begin{cases}
    f(x) :: \mathsf{map}(f, \mathsf{filter}(p \circ f, \mathit{xs}))
      & \text{if } (p \circ f)(x) \\
    \mathsf{map}(f, \mathsf{filter}(p \circ f, \mathit{xs}))
      & \text{otherwise}
  \end{cases} \quad [\text{IH}] \\
  &= \mathsf{map}(f, \mathsf{filter}(p \circ f, x::\mathit{xs}))
\end{align*}


\section{Example 12: Insertion Sort is Stable}

\textbf{Strategy}: Loop invariant.

This proves a property that \emph{cannot} be proved by sampling or Z3:
stability (equal elements maintain their relative order).

\textbf{Invariant}: $\mathit{sorted}[0..i]$ is a stable permutation
of $\mathit{xs}[0..i]$.

\textbf{Preservation}: inserting $\mathit{xs}[i]$ into the sorted
prefix at the correct position preserves stability because insertion
happens \emph{after} equal elements (not before), so equal elements
maintain their original order.

This is a deeply structural argument that requires understanding the
insertion algorithm's behavior on equal elements---exactly the kind
of reasoning that proof terms enable.


\section{Proof Complexity Analysis}

\begin{longtable}{lccl}
\toprule
\textbf{Example} & \textbf{Depth} & \textbf{Size} & \textbf{Strategy} \\
\midrule
Factorial rec/iter & 3 & 4 & NatInduction \\
Merge/Quick sort & 2 & 3 & AbsRef \\
Edit distance DP/memo & 2 & 4 & Simulation \\
BFS/Union-Find & 2 & 3 & AbsRef \\
Dijkstra/Bellman-Ford & 2 & 4 & Simulation \\
Trial div/sieve & 2 & 3 & AbsRef \\
Powerset rec/bitmask & 2 & 3 & CommDiag \\
LCS DP/Hunt-Szymanski & 2 & 4 & Simulation \\
RBT/AVL insert & 2 & 3 & AbsRef \\
NFA/DFA matching & 2 & 4 & Simulation \\
Map-filter commute & 3 & 4 & ListInduction \\
Stable insertion sort & 2 & 5 & LoopInvariant \\
Gauss sum & 3 & 4 & NatInduction \\
reverse(reverse) & 3 & 4 & ListInduction \\
$2 \cdot 3 + 4 = 10$ & 1 & 1 & ArithSimp \\
$1 \cdot x + 0 = x$ & 5 & 5 & Trans+Cong+Axiom \\
\bottomrule
\end{longtable}

\begin{remark}
Proof depth is typically 2--5.  The most complex proofs (equational
chains like $1 \cdot x + 0 = x$) have depth 5 because they chain
multiple rewrite steps.  Algorithmic equivalences (simulation, abstraction
refinement) have low depth (2--3) because the complexity is in the
\emph{relation} or \emph{specification}, not the proof structure.
\end{remark}
